\section*{Bab \ref{ch:b2:mass-spec-atomic} --- Spektrometri Massa dan Spektroskopi Atom}

\begin{solution}{exo:b2:mass-spec-atomic:1}
78; $400/2 = 200$; $194 + 1 = 195$, $[\mathrm M + \mathrm H]^+$.
\end{solution}

\begin{solution}{exo:b2:mass-spec-atomic:2}
Massa ganjil: jumlah atom nitrogen ganjil, satu untuk molekul sekecil itu.
\ce{C3H7NO} (propanamida, 73) atau \ce{C4H11N} (butan-1-amina, 73).
\end{solution}

\begin{solution}{exo:b2:mass-spec-atomic:3}
Hampir sama: satu brom. 3 : 1: satu klor.
\end{solution}

\begin{solution}{exo:b2:mass-spec-atomic:4}
Kuning: natrium, \qty{589.0}{nm} dan \qty{589.6}{nm}. Merah tua: litium, \qty{670.8}{nm}.
% ledger: asd:NaD2, asd:NaD1, asd:Li671
\end{solution}

\begin{solution}{exo:b2:mass-spec-atomic:5}
$n_C \approx 6.7/1.11 \approx 6$.
\end{solution}

\begin{solution}{exo:b2:mass-spec-atomic:6}
\ce{CO}: 27.9949; \ce{N2}: 28.0061; \ce{C2H4}: 28.0313 u. Daya pisah: \ce{CO}/\ce{N2} $28/0.0112
\approx 2.5 \times 10^3$; \ce{N2}/\ce{C2H4} $28/0.0252 \approx 1.1 \times 10^3$; \ce{CO}/\ce{C2H4}
$28/0.0364 \approx 7.7 \times 10^2$.
% ledger: mass:C12, mass:O16, mass:N14, mass:H1
\end{solution}

\begin{solution}{exo:b2:mass-spec-atomic:7}
86: \omterm{def:b2:mass-spec-atomic:molecular-ion}{ion molekul} \ce{C5H10O}. 57: \ce{C2H5CO+}, \omterm{def:b2:aromatic-substitution:friedel-crafts}{ion asilium} hasil
pemutusan-$\alpha$ (kehilangan radikal etil, 29). 29: \ce{C2H5+}.
% ledger: ms:pentan-3-one
\end{solution}

\begin{solution}{exo:b2:mass-spec-atomic:8}
86: $\mathrm M^{+\bullet}$. 71: $\mathrm M - \ce{CH3}$, \omterm{def:b2:aromatic-substitution:friedel-crafts}{ion asilium}
\ce{C3H7CO+}. 43: \ce{CH3CO+}
(kehilangan radikal propil). 58: \omterm{def:b2:mass-spec-atomic:fragmentation}{penataan ulang McLafferty}, melalui
hidrogen pada karbon
$\gamma$ rantai propil, dengan kehilangan etena. Dalam pentan-3-on,
gugus-gugus etil tidak mempunyai karbon $\gamma$: tidak ada ion McLafferty;
puncak kecilnya pada 58 adalah puncak isotop \ce{^{13}C} dari ion pada 57
(tiga karbon, sekitar 3.3~\%).
% ledger: ms:pentan-2-one, ms:pentan-3-one
\end{solution}

\begin{solution}{exo:b2:mass-spec-atomic:9}
Garis yang melalui titik asal: kemiringan $\sum c_iA_i/\sum c_i^2 = 1.553/30.0 \approx \qty{0.0518}{L/mg}$. Sampel:
$0.128/0.0518 \approx \qty{2.47}{mg/L}$ tembaga.
\end{solution}

\begin{solution}{exo:b2:mass-spec-atomic:10}
$t = L\sqrt{m/(2eU)} = 1.00\sqrt{1000 \times 1.6605 \times 10^{-27}/(2 \times 1.6022 \times 10^{-19}
\times 2.00 \times 10^4)} \approx \qty{16.1}{\micro s}$. Untuk 1001 u, $t$ bertambah dengan faktor
$\sqrt{1.001} \approx 1 + 0.0005$: $\Delta t \approx \qty{8.0}{ns}$.
% ledger: const:u, const:e
\end{solution}

\begin{solution}{exo:b2:mass-spec-atomic:11}
Dua klor: $(0.758 + 0.242x)^2$ memberikan $0.575$, $0.367$, $0.0586$: $m/z$
84, 86, dan 88 dengan rasio
100 : 64 : 10.
\end{solution}

\begin{solution}{exo:b2:mass-spec-atomic:12}
\ce{^{40}Ar^{16}O+}: $39.9624 + 15.9949 = 55.9573$; \ce{^{56}Fe}: 55.9349; $R = 56/0.0224 \approx 2.5 \times
10^3$. Cara lain: ukurlah isotop besi yang lain, \ce{^{57}Fe}, atau
singkirkan \ce{ArO+} dalam sel tumbukan sebelum penganalisis.
% ledger: mass:Ar40, mass:Fe56, mass:O16
\end{solution}

\begin{solution}{pb:b2:mass-spec-atomic:1}
\textbf{1.} 156, 158, dan 160 (bersama 157 dan 159): molekul itu berada
dalam beberapa bentuk isotopik.
\textbf{2.} Tiga puncak yang berjarak dua satuan, dengan puncak tengah yang
terbesar: satu klor dan satu brom.
\textbf{3.} 156.
\textbf{4.} Massa genap: tidak ada nitrogen, atau jumlahnya genap.
\textbf{5.} $156 - 79 - 35 = 42$: \ce{C3H6}.
\textbf{6.} \ce{C3H6BrCl}; derajat ketidakjenuhan $(2 \times 3 + 2 - 6 - 2)/2 = 0$:
tidak ada cincin, tidak ada ikatan rangkap.
\textbf{7.} $0.5/14.5 \approx 3.4~\%$, $3.4/1.11 \approx 3$ karbon.
\textbf{8.} Puncak pada 157 diberikan sampai 0.1 pada skala yang nilainya
0.5: ketidakpastian 20~\%.
\textbf{9.} \ce{C3H6^{81}Br^{35}Cl} dan \ce{C3H6^{79}Br^{37}Cl}.
\textbf{10.} $3 \times 12 + 6 \times 1.007825 + 78.918338 + 34.968853 \approx 155.9341$.
\textbf{11.} $156/(156.151 - 155.934) \approx 7.2 \times 10^2$.
\textbf{12.} Ya: dua hidrogen lebih sedikit, 154.
\textbf{13.} Sebuah atom brom (79 atau 81) hilang: \ce{C3H6Cl+}, 77 dengan
\ce{^{35}Cl}, 79 dengan \ce{^{37}Cl},
3 : 1.
\textbf{14.} Kehilangan HBr: kation radikal $\mathrm{C_3H_5Cl^{+\bullet}}$, bermassa genap.
\textbf{15.} \ce{C3H5+}, kation alil, setelah kehilangan kedua halogen (Br, lalu HCl).
\textbf{16.} \ce{CH2Cl+} (49 dan 51, 3 : 1).
\textbf{17.} \ce{CH2Br+} (93 dan 95) dan \ce{C2H4Br+} (107 dan 109),
masing-masing 1 : 1: brom berada pada salah satu ujung rantai, ion kedua
berasal dari kehilangan \ce{CH2Cl}.
\textbf{18.} Ikatan \ce{C-Br} lebih lemah daripada ikatan \ce{C-Cl}, dan
\ce{Br} adalah radikal pergi yang lebih baik.
\textbf{19.} Tidak: tidak ada puncak pada 127--131; \ce{CH2Cl+} dan
\ce{CH2Br+} menunjukkan bahwa setiap halogen berada pada \ce{CH2} ujung.
\textbf{20.} \ce{BrCH2CH2CH2Cl}, 1-bromo-3-kloropropana.
\textbf{21.} 156/158/160: M; 77/79: M $-$ Br; 76: M $-$ HBr; 107/109: M $-$ \ce{CH2Cl};
93/95: \ce{CH2Br+}; 49/51: \ce{CH2Cl+}; 41: \ce{C3H5+}; 39, 27: \ce{C3H3+}, \ce{C2H3+}.
\textbf{22.} Tidak: tidak ada gugus karbonil.
\textbf{23.} 1-Bromo-2-kloropropana, \ce{CH3CHClCH2Br}: kehilangan \ce{CH2Br}
menghasilkan \ce{CH3CHCl+} pada 63/65.
\textbf{24.} \ce{C3H6^{81}Br^{37}Cl}.
\textbf{25.} $M$: $0.758 \times 0.5069 = 0.384$; $M+2$: $0.758 \times 0.4931 + 0.242 \times 0.5069 =
0.496$; $M+4$: $0.242 \times 0.4931 = 0.119$. Rasio $\boldsymbol{\approx 3.2 : 4.2 : 1}$,
sekitar 3 : 4 : 1; spektrumnya memberikan $14.5 : 18.6 : 4.4 = 3.3 : 4.2 : 1$.
\end{solution}
